% =====================================================================
%  AyurLLM — CVPR-style two-column paper
%  Compile:  pdflatex ayurllm_paper.tex   (twice, for cross-references)
%
%  If you have the official CVPR style files, delete the GEOMETRY block
%  below and use \usepackage{cvpr} (or \usepackage[review]{cvpr}) instead.
%  The geometry here reproduces the CVPR text block exactly.
% =====================================================================
\documentclass[10pt,twocolumn,letterpaper]{article}

\usepackage[T1]{fontenc}
\usepackage[utf8]{inputenc}
\usepackage{times}
\usepackage{graphicx}
\usepackage{amsmath,amssymb}
\usepackage{booktabs}
\usepackage{multirow}
\usepackage{array}
\usepackage{xcolor}
\usepackage{microtype}
\usepackage{caption}
\usepackage[colorlinks=true,linkcolor=black,citecolor=black,urlcolor=black]{hyperref}

% ---------------- GEOMETRY: CVPR text block ----------------
\usepackage[letterpaper,
            top=1in, bottom=1.125in,
            left=0.84375in, right=0.84375in]{geometry}
\setlength{\columnsep}{0.3125in}
% -----------------------------------------------------------

\captionsetup{font=small,labelfont=bf,skip=4pt}
\setlength{\parskip}{0pt}
\setlength{\abovedisplayskip}{4pt}
\setlength{\belowdisplayskip}{4pt}

% ---- the final OCR engine is set in ONE place; change it here ----
% If "LiFT" is in fact Liquid AI's LFM2-VL, change the line below to
% \newcommand{\ocrfinal}{LFM2-VL (GGUF)}  and the whole paper updates.
\newcommand{\ocrfinal}{LiFT (GGUF)}

% ---- placeholders for values not yet supplied -----------------------
% \draftgapstrue   shows every unfinished section as a boxed to-do  (working copy)
% \draftgapsfalse  hides them, leaving only the finished text       (copy to share)
\newif\ifdraftgaps
\draftgapstrue

\newcommand{\TBD}[1]{\ifdraftgaps\textcolor{red}{\textbf{[#1]}}\else\rule{1.2em}{0pt}\fi}
\newcommand{\gap}[2]{%
  \ifdraftgaps
    \par\medskip\noindent\fbox{\parbox{0.95\columnwidth}{%
    \small\textbf{[ TO BE COMPLETED --- #1 ]}\\[2pt]\emph{#2}}}\par\medskip
  \fi}

\title{\vspace{-1.2em}\textbf{\fontsize{14}{16}\selectfont
AyurLLM: A Verified Shloka-Level Corpus and a Traceable Instruction
Benchmark from Eight Classical Ayurvedic Treatises}\vspace{-0.4em}}

\author{Anonymous CVPR submission\\
Paper ID \TBD{ID}}
\date{}                     % CVPR papers carry no date line

\begin{document}
\maketitle

% =====================================================================
\begin{abstract}
Classical Ayurvedic knowledge is distributed as printed treatises in which
Devanagari verse, IAST transliteration and prose commentary in English or
Hindi are interleaved on the page, and in which the same logical unit --- a
\emph{shloka} and its gloss --- is typeset differently in every edition. Such
material is not searchable by symptom, and it is not directly usable as
training data. We present AyurLLM, an end-to-end pipeline that converts eight
such treatises into a verified, schema-normalised corpus and into an
instruction benchmark whose every answer remains traceable to the printed
page. We report four contributions. First, an OCR selection protocol over six
extraction strategies, in which layout-aware and vision-language extraction
are retained and general-purpose detector--recogniser pipelines are rejected
on Devanagari conjunct handling rather than on aggregate accuracy. Second, a
single chunk schema that every book-specific converter targets, yielding
5{,}838 verified chunks and 2{,}708{,}568 tokens under a production
tokenizer, of which 51.87\% are Sanskrit, 45.67\% English and 2.47\% Hindi.
Third, a per-volume failure analysis that attributes each gap in coverage to a
specific, named extraction defect --- Hindi-only commentary, degraded
Devanagari, absent verse fields, a single catastrophic recognition failure,
and cross-volume identifier collision --- rather than to source quality in
aggregate. Fourth, a fact-centric synthesis procedure that reduces each chunk
to discrete facts and renders each fact into every examination format it
supports, producing 80{,}000 questions over six formats from 35{,}016 facts
with no generative model in the answer path, and therefore no possibility of
a hallucinated answer. We additionally report a paired model-selection study
over four instruction-tuned candidates that identifies both a preferred
fine-tuning base and two question formats on which no candidate exceeds
chance.
\end{abstract}

% =====================================================================
\section{Introduction}

Ayurveda is a codified medical tradition whose primary sources are a small,
stable set of Sanskrit treatises. A practising physician reasons from a
symptom, a drug or a disease name; the literature is organised by verse
number within chapter within section. The two do not meet. A physician who
wishes to know what the tradition says about a presenting symptom must
already know which verse to look up, which is precisely the knowledge the
lookup was meant to supply.

Making these texts machine-searchable is not primarily a modelling problem.
It is a problem of getting the text off the page correctly, and of deciding
what the unit of retrieval should be. Both are harder than they appear.
Devanagari is a complex script in which stacked conjunct consonants, vowel
signs positioned above, below and beside their base glyph, and the connecting
headline interact in ways that general-purpose recognisers handle
poorly~\cite{devaocr}. Sanskrit compounds these difficulties: the same
printed page may carry verse in Devanagari, its romanisation in IAST, and a
commentary that is itself in Sanskrit, Hindi or English. Editions disagree on
all of it.

The unit question is equally consequential. A verse without its commentary is
often uninterpretable; a commentary block without its verse cannot be cited.
We therefore treat the pair --- verse together with its gloss and its
position in the text --- as the atomic record, and require every downstream
stage to operate on that record rather than on undifferentiated text.

A further difficulty arises once the corpus exists. Instruction tuning
requires supervised examples that a corpus of primary sources does not
natively contain, and the standard remedy is to generate them with a language
model. In a medical domain this is a poor trade. A teacher model that
hallucinates a detail does not produce a noisy label; it produces a confident
false label that becomes the student's ground truth~\cite{attrground}. We
consider this unacceptable for a corpus whose purpose is to tell a clinician
what a classical text actually says, and we therefore constrain synthesis so
that every answer is lifted verbatim from a source chunk.

This paper makes the following contributions.

\begin{itemize}\itemsep2pt
\item \textbf{An OCR selection protocol for Indic scriptural material}
(Sec.~\ref{sec:ocr}). We evaluate six extraction strategies and report which
were accepted and which rejected, with the failure mode that decided each.
The deciding criterion is conjunct and matra fidelity on verse, not aggregate
page accuracy.

\item \textbf{A single chunk schema and its enforcement}
(Sec.~\ref{sec:schema}). Every book-specific converter targets one contract,
so that irregularity is absorbed at extraction time rather than propagated.
We also report a cross-volume identifier collision that per-file validation
cannot detect.

\item \textbf{A per-volume failure analysis} (Sec.~\ref{sec:failure}).
We attribute each coverage gap to a specific defect with the evidence that
identifies it, and show that token fertility functions as a diagnostic for
recognition quality.

\item \textbf{A fact-centric, hallucination-free instruction corpus}
(Sec.~\ref{sec:synthesis}) of 80{,}000 questions over six formats, together
with a paired model-selection study (Sec.~\ref{sec:selection}) that
identifies a fine-tuning base and quantifies where current models fail.
\end{itemize}

% =====================================================================
\section{Related Work}
\label{sec:related}

\subsection{Document conversion and OCR}

Modern document conversion has moved away from the classical
detect-then-recognise cascade toward systems that model page structure
explicitly. PP-OCR~\cite{ppocr} is representative of the cascade approach and
optimises it aggressively for deployment, reaching a 3.5\,MB recognition
model for 6{,}622 Chinese characters through a sequence of size-reduction
strategies; PP-OCRv2~\cite{ppocrv2} extends this with a further set of
training refinements. This lineage is fast and well-engineered, but its
segmentation assumptions are inherited from scripts whose glyphs do not stack
vertically.

Docling~\cite{docling,doclingtr} takes the structural view, parsing documents
into a unified representation using dedicated layout-analysis and
table-structure models, and is explicitly designed to run within a small
resource budget on commodity hardware. Surya~\cite{surya} folds layout
analysis, reading-order detection, table recognition and text recognition
into a single vision--language model covering over ninety languages, and
Marker~\cite{marker} uses Surya as its recognition backend while targeting
whole-document conversion to structured Markdown or JSON. The trend these
represent --- from pipeline to end-to-end vision--language model --- is
directly relevant to scriptural material, where the correct reading of a
glyph often depends on layout context that a cropped line image discards.

Devanagari recognition has its own literature. Surveys of the
area~\cite{devaocr} identify conjunct consonants, dependent vowel signs and
the connecting headline as the persistent failure points, and recent
stress-testing of vision--language OCR systems on Devanagari~\cite{devavlm}
confirms that these remain unsolved even for models that perform well on
Latin and CJK text. Work specific to Sanskrit manuscripts~\cite{sanskritdl}
notes the additional problem that the effective character inventory is
unusually large once compound characters are counted. Our own findings in
Sec.~\ref{sec:failure} are consistent with this: the volume in our corpus
with the worst recognition quality fails precisely by losing conjuncts and
detaching matras.

\subsection{Domain and language adaptation}

Adapting general models to low-resource languages and to specialist domains
is typically approached through continued pre-training, and empirical studies
of this for Indic languages~\cite{cpt-indic} examine the interaction between
corpus size and vocabulary extension. Surveys of South Asian language
resources~\cite{bhaasha} document how thin the available data remains for
most of these languages, which is the condition our corpus construction
operates under.

The most direct precedent for this work is AyurParam~\cite{ayurparam}, a
bilingual model specialised to Ayurveda, evaluated on the
BhashaBench~\cite{bhashabench} suite. That work and ours differ in the unit
of grounding: AyurParam is trained on page-level question--answer pairs
generated across a large book collection, whereas we ground at the level of
the individual verse and its commentary, with an explicit provenance chain
back to the printed page. We treat AyurParam as the reference point against
which our fine-tuned model must be measured (Sec.~\ref{sec:ayurparam}).
Broader Indic evaluation suites~\cite{parambench,indicparam} provide
complementary coverage at graduate level and across low-resource Indic
languages.

\subsection{Synthetic instruction data}

Bootstrapping instruction data from a model's own generations is now standard
practice, and subsequent work has concentrated on raising diversity and
difficulty while controlling quality. Attributed grounding~\cite{attrground}
anchors each generated instruction to retrieved source material precisely to
limit fabrication, and answer-consistency filtering~\cite{cotself} curates
generated data by checking that independently sampled reasoning paths agree.
Seed-free frameworks~\cite{seedfree} address the case where no in-language
instruction seed set exists, which is the situation for Sanskrit.

Our position is more conservative than any of these. Rather than generating
answers and then filtering, we do not generate answers at all: every answer
in our instruction corpus is a verbatim span of a source chunk, and the
generative step is confined to selecting which span to ask about and in which
examination format. This forfeits the linguistic variety that a teacher model
would supply --- a limitation we quantify in Sec.~\ref{sec:limitations} ---
in exchange for the guarantee that no answer can be false with respect to the
source.

% =====================================================================
\section{Corpus Construction}
\label{sec:corpus}

\subsection{Source material}

The corpus is drawn from eight printed volumes selected in consultation with
practising physicians, spanning the canonical divisions of the literature:
the \emph{Brihattrayi} (Charaka Samhita, Sushruta Samhita, Ashtanga Hridaya)
and the \emph{Laghutrayi} (Sharangadhara Samhita, Madhava Nidana,
Bhavaprakasha Nighantu), together with a two-volume edition of Madhava
Nidanam carrying the Madhukosa commentary.

The eight volumes differ in every respect that matters to an extraction
pipeline: in script mix, in whether transliteration is printed at all, in
whether the commentary is English or Hindi, in scan quality, and in whether
verse and commentary occupy separate columns or a single measure. No single
extraction configuration succeeds across all eight. This is the premise of
the engine selection described next, and the cause of the failures analysed
in Sec.~\ref{sec:failure}.

\subsection{OCR engine selection}
\label{sec:ocr}

Extraction was treated as a per-volume engineering problem. Six strategies
were evaluated. The deciding criterion was not aggregate page accuracy but
fidelity on Devanagari verse specifically: whether conjunct consonants
survive intact and whether dependent vowel signs remain attached to the
correct base glyph. A page can score well on character accuracy while
rendering every verse on it unusable, because the commentary is long and the
verse is short.

\paragraph{Accepted.}
\textbf{\ocrfinal{}} was adopted as the final configuration. Quantisation was
a hard requirement rather than an optimisation: extraction ran on a single
local accelerator with approximately 6\,GB of usable memory, which precludes
full-precision inference over multi-hundred-page volumes. GGUF-quantised
vision--language models are explicitly built for this
regime~\cite{lfm2vl}. \textbf{Surya}~\cite{surya} was accepted for
verse-dominant pages, where its unified treatment of layout and recognition
preserved reading order across interleaved verse and commentary.
\textbf{Marker}~\cite{marker} was accepted where the printed page carried a
stable heading hierarchy that could be mapped onto the schema's chapter and
section fields. \textbf{Docling}~\cite{docling} was accepted where verse and
commentary occupied physically distinct page regions, its layout model
being able to keep them apart. \textbf{Explicit bounding-box extraction} was
retained as a fallback for columns that the general-purpose converters
merged.

\paragraph{Rejected.}
\textbf{PaddleOCR / PP-OCR}~\cite{ppocr,ppocrv2} was rejected. Its
detect-then-recognise cascade segments text into line and character regions
before recognition, and this assumption is a poor fit for stacked conjuncts,
where the correct unit is neither a character box nor a clean horizontal
line. The engine is otherwise the most deployment-efficient of those
evaluated, and the rejection is specific to the script rather than a
judgement on the system. \textbf{Markitdown} was rejected as a primary
extractor: it is a format converter rather than a recogniser and produced
usable output only where the source PDF already carried a text layer, which
was not the case for the scanned volumes.

\begin{table}[t]
\centering
\small
\setlength{\tabcolsep}{4pt}
\begin{tabular}{@{}llc@{}}
\toprule
\textbf{Engine} & \textbf{Role} & \textbf{Decision} \\
\midrule
\ocrfinal{}          & Final extractor, all volumes & Accepted \\
Surya~\cite{surya}   & Verse-dominant pages         & Accepted \\
Marker~\cite{marker} & Stable heading hierarchy     & Accepted \\
Docling~\cite{docling} & Separated page regions     & Accepted \\
Bounding-box         & Merged-column fallback       & Accepted \\
PP-OCR~\cite{ppocr}  & Conjunct segmentation        & Rejected \\
Markitdown           & No text layer present        & Rejected \\
\bottomrule
\end{tabular}
\caption{Extraction strategies evaluated, with the role each was assigned or
the failure mode that excluded it.}
\label{tab:ocr-decision}
\end{table}

\begin{table}[t]
\centering
\small
\setlength{\tabcolsep}{3pt}
\begin{tabular}{@{}lcccc@{}}
\toprule
\textbf{Engine} & \textbf{CER} & \textbf{WER} & \textbf{Conj.\ err.} & \textbf{p/min} \\
\midrule
\ocrfinal{}   & \TBD{--} & \TBD{--} & \TBD{--} & \TBD{--} \\
Surya         & \TBD{--} & \TBD{--} & \TBD{--} & \TBD{--} \\
Marker        & \TBD{--} & \TBD{--} & \TBD{--} & \TBD{--} \\
Docling       & \TBD{--} & \TBD{--} & \TBD{--} & \TBD{--} \\
Bounding-box  & \TBD{--} & \TBD{--} & \TBD{--} & \TBD{--} \\
PP-OCR        & \TBD{--} & \TBD{--} & \TBD{--} & \TBD{--} \\
Markitdown    & \TBD{--} & \TBD{--} & \TBD{--} & \TBD{--} \\
\bottomrule
\end{tabular}
\caption{Measured extraction quality on the held-out transcription set.
Character error rate, word error rate, conjunct-specific error rate and
throughput in pages per minute.}
\label{tab:ocr-stats}
\end{table}

\gap{OCR measurements for Table~\ref{tab:ocr-stats}}{%
Supply per-engine CER, WER, conjunct-specific error rate and pages per
minute, plus peak memory if available. Also state the ground-truth protocol:
how many pages were manually transcribed, from which volumes, and how they
were sampled. Until these are filled the accept/reject decisions in
Sec.~\ref{sec:ocr} rest on qualitative failure modes only, which a reviewer
will challenge.}

Extraction was made resumable. Process isolation and per-page checkpointing
ensured that a crash partway through a volume did not discard completed
pages, which within a 6\,GB budget is a routine occurrence rather than an
exceptional one.

\subsection{The chunk schema}
\label{sec:schema}

Every book-specific converter targets one common contract. This is the
central design decision of the corpus: converters absorb the irregularity of
the printed page so that all downstream stages --- tokenization, embedding,
question synthesis, evaluation --- operate on a single uniform record type.
The base schema is: \texttt{chunk\_id}; \texttt{sthana},
\texttt{chapter\_number} and \texttt{chapter\_title} locating the record in
the volume; \texttt{verse\_id} as printed; \texttt{topic\_heading}, the
editorial heading under which the verse appears; \texttt{sanskrit\_text},
the Devanagari verse; \texttt{transliteration}, the IAST romanisation where
the edition prints one; \texttt{commentary}, the prose gloss retained
verbatim; and \texttt{categories\_found} with \texttt{primary\_category},
domain keyword tags used for filtering and downstream domain assignment.

Three volumes required schema extension rather than being forced into the
base contract. Sharangadhara Samhita carries
\texttt{commentary\_english}, because its printed commentary is Hindi and an
English rendering had to be produced separately. Madhava Nidana (Laghutrayi)
carries \texttt{sanskrit\_commentary}, \texttt{source\_citation} and
\texttt{record\_type}, because its Madhukosa apparatus quotes other treatises
by name and the citation must be preserved. Charaka Samhita carries
\texttt{pada}, marking the quarter-division of its Chikitsasthana chapters.
Retaining these as named fields rather than flattening them into
\texttt{commentary} is what makes the failure analysis of
Sec.~\ref{sec:failure} possible at all.

\subsection{Identifier integrity}

\texttt{chunk\_id} is unique within a volume but not across the corpus.
Madhava Nidanam Part 1 and Part 2 were converted independently and both
number their records \texttt{MN1}, \texttt{MN2}, and so on; 102 identifiers
collide between the two volumes. Any join keyed on \texttt{chunk\_id} alone
therefore silently double-counts. All analysis in this paper keys on the
composite \texttt{(source\_book, chunk\_id)}. We report this because it is
invisible to per-volume validation --- each file is internally consistent ---
and surfaces only when the corpus is treated as a whole.

% =====================================================================
\section{Corpus Characterisation}
\label{sec:characterisation}

\subsection{Tokenization}

The corpus was tokenized with the \texttt{google/gemma-3-4b-it} tokenizer,
chosen because it belongs to a model family in the deployment size class
under consideration. Measuring token counts under a tokenizer other than the
one that will be used would misstate every context-window and cost estimate
that follows.

The corpus totals 2{,}708{,}568 tokens: 1{,}404{,}844 Sanskrit (51.87\%),
1{,}236{,}881 English (45.67\%) and 66{,}843 Hindi (2.47\%). Devanagari-only
content accounts for 801{,}705 tokens across 5{,}439 chunks, and
transliteration-only content for 572{,}022 tokens across 2{,}950 chunks ---
only about half the corpus carries printed romanisation at all.

\subsection{Token fertility as a quality diagnostic}

Fertility, tokens emitted per whitespace-delimited word, is 3.61 for Sanskrit
against 1.69 for English, a ratio of 2.14: Sanskrit is more than twice as
expensive per word under this tokenizer as the English commentary explaining
it. Fertility is not uniform, ranging from 1.85 in Ashtanga Hridaya to 4.70
in Sushruta Samhita (Table~\ref{tab:corpus}).

The low Ashtanga figure is not evidence of simpler Sanskrit. It is a symptom.
Lost conjuncts and detached vowel signs fragment words into shorter,
tokenizer-friendly pieces, so degraded recognition \emph{lowers} apparent
fertility. Fertility therefore doubles as a diagnostic for extraction
quality, and an anomalously low value is a signal to inspect the volume
rather than a property of its language.

\begin{table}[t]
\centering
\small
\setlength{\tabcolsep}{4pt}
\begin{tabular}{@{}lrrr@{}}
\toprule
\textbf{Volume} & \textbf{Chunks} & \textbf{Skt tokens} & \textbf{Fert.} \\
\midrule
Charaka Samhita          & 1{,}497 & 379{,}568 & 3.41 \\
Bhavaprakasha Nighantu   & 1{,}052 & 81{,}372  & 3.69 \\
Ashtanga Hridaya         & 913     & 38{,}899  & 1.85 \\
Madhava Nidanam P1       & 728     & 223{,}062 & 4.24 \\
Sharangadhara Samhita    & 593     & 61{,}766  & 4.43 \\
Sushruta Samhita         & 418     & 58{,}848  & 4.70 \\
Madhava Nidana (Laghu.)  & 403     & 259{,}833 & 4.09 \\
Madhava Nidanam P2       & 234     & 301{,}496 & 4.65 \\
\midrule
\textbf{Total}           & \textbf{5{,}838} & \textbf{1{,}404{,}844} & \textbf{3.61} \\
\bottomrule
\end{tabular}
\caption{Per-volume chunk counts, Sanskrit token totals and Sanskrit token
fertility. The anomalously low Ashtanga fertility is a recognition artefact,
not a linguistic property (Sec.~\ref{sec:failure}).}
\label{tab:corpus}
\end{table}

\subsection{Chunk length and context limits}

Chunk length is heavily skewed. The median is comfortably inside any modern
context window; the upper tail is not. 21.62\% of chunks (1{,}262) exceed
512 tokens, 8.56\% (500) exceed 1{,}024, 2.93\% (171) exceed 2{,}048, 0.87\%
(51) exceed 4{,}096 and 0.15\% (9) exceed 8{,}192.

The skew is also uneven across volumes. Madhava Nidanam Part 2 averages
1{,}288 Sanskrit tokens per chunk against 43 for Ashtanga Hridaya, a
thirtyfold difference driven by editorial convention rather than content
density: Part 2 emits one chunk per commentary block, Ashtanga one per verse.
Operationally, chunks above the embedding model's context limit must be
sub-chunked before ingestion or their tails are silently truncated. We treat
the 512-token threshold as binding, since it governs roughly one chunk in
five.

\subsection{Script balance}

The ratio of English to Sanskrit tokens per chunk separates the corpus into
two classes. In Bhavaprakasha Nighantu the mean ratio is 3.85: it is a
\emph{nighantu}, a materia medica whose substance is prose description of
drug properties, with verse serving as an index. In Sushruta Samhita the
ratio is 0.48 and in Madhava Nidanam Part 2 it is 0.46, where verse dominates
and commentary is terse. Any downstream sampling that treats chunks as
interchangeable will over-represent pharmacological prose and
under-represent surgical and diagnostic verse. Sec.~\ref{sec:synthesis}
returns to this.

% =====================================================================
\section{Per-Volume Failure Analysis}
\label{sec:failure}

Aggregate quality statistics conceal the fact that each volume fails
differently. We report the dominant defect per volume with the evidence that
identifies it. In every case the defect is an extraction or editorial
artefact, not absent source content; this was confirmed by returning to the
printed page.

\paragraph{Sharangadhara Samhita --- Hindi-only commentary.}
350 of 593 chunks carry no English commentary. The printed commentary is
Hindi, sharing the Devanagari script with the Sanskrit verse and offering no
English text to anchor language identification on. The corpus records 560
chunks with Hindi content totalling 66{,}843 tokens: the entire Hindi mass of
the corpus originates in this single volume. Character-level translation
could not reliably separate Hindi commentary from Sanskrit verse and a
meaning-level pass was required. Only 42.2\% of this volume's chunks carry
both a verse and usable English commentary, making it the weakest volume for
any English-language downstream task.

\paragraph{Ashtanga Hridaya --- degraded Devanagari.}
This volume shows the corpus minimum on both Sanskrit fertility (1.85) and
mean Sanskrit chunk length (43.51 tokens). Inspection of affected chunks
identifies the cause: conjunct consonants lost and vowel signs detached from
their base characters during recognition, fragmenting words into isolated
glyphs. The text remains legible enough to pass a naive non-empty check while
being unusable for any task requiring the verse to be read. This defect class
is invisible to field-coverage metrics, and motivated the explicit Devanagari
well-formedness test described in Sec.~\ref{sec:controls}.

\paragraph{Bhavaprakasha Nighantu --- absent verse field.}
370 of 1{,}052 chunks carry no Sanskrit text; the shortest chunks in the
corpus, at three to four tokens, are from this volume. These are drug
monograph entries whose printed page carries prose description without an
accompanying verse. The records are legitimate content but cannot support any
verse-grounded task.

\paragraph{Charaka Samhita --- catastrophic single-chunk failure.}
One chunk, \texttt{CS1543}, contains 116{,}187 combined tokens: nine times
the next largest in the corpus and roughly 4\% of total corpus token mass in
a single record. Its content is not Ayurvedic text but the volume's back
matter, rendered by the recognition pass as a long run of fragmentary glyphs
and page numbers. A single record of this kind distorts every mean it enters,
raising this volume's English standard deviation to 2{,}558 tokens against a
median of 110.

\paragraph{Madhava Nidanam Parts 1 and 2 --- identifier collision.}
The 102 shared identifiers described in Sec.~\ref{sec:schema}. Part 1
additionally shows 130 of 728 chunks with no English commentary.

\paragraph{Remaining volumes.}
Sushruta Samhita and Madhava Nidana (Laghutrayi) show no dominant structural
defect. Sushruta is complete on both fields across all 418 chunks;
Laghutrayi carries 26 chunks without English commentary out of 403.
Sushruta's terse commentary --- 66.35 English tokens per chunk, the corpus
minimum --- is an editorial property of the edition rather than a failure.

% =====================================================================
\section{Embedding and Retrieval}
\label{sec:embedding}

Embedding was performed with Harrier OSS over the schema-normalised chunks.
The retrieval layer is a hybrid of lexical and dense search, so that a
physician may query either by an exact Sanskrit term or by a plain-language
description of a symptom, and is served through a search backend with a web
interface, a REST API and a machine-facing endpoint for programmatic access
by assistant models.

\gap{Embedding and index configuration}{%
Needed: embedding model revision and output dimension; whether sub-chunking
was applied to the 1{,}262 chunks above 512 tokens and with what stride and
overlap; which schema fields were concatenated to form the embedded text; the
index configuration including searchable and filterable attributes; and the
lexical/dense weighting at query time. If retrieval quality was measured,
supply the query set and metric.}

% =====================================================================
\section{Instruction Data Synthesis}
\label{sec:synthesis}

Fine-tuning requires supervised examples the corpus does not natively
contain. We synthesise them under the constraint motivated in
Sec.~\ref{sec:related}: every answer must be extracted verbatim from a source
chunk and remain traceable to it. No generative model participates in
producing answer content, so no answer can be hallucinated, and any item can
be audited by returning to its \texttt{(source\_book, chunk\_id)} pair.

\subsection{Fact-centric generation}

Rather than generating questions directly from text, each chunk is first
reduced to a set of discrete facts, each carrying a stable identifier. A fact
may be the subject a verse addresses, a statement made in the commentary, a
technical term the passage turns on, an explicit enumeration or count, a
synonym set, or the location of a teaching within the volume. Each fact is
then rendered into every examination format it can support.

This is what allows one answer to be tested as a multiple-choice item, a
cloze, an assertion--reason pair and a multiple-selection item without any two
question stems being textually identical, and it makes semantic
de-duplication possible by grouping on the fact identifier. From 5{,}838
chunks we extract 35{,}016 distinct facts and render 80{,}000 questions
across six formats, covering 93.9\% of chunks; 21{,}635 facts (61.8\%) are
rendered in more than one format. The residual 6.1\% of uncovered chunks are
those whose only content is Hindi-only commentary or degraded Devanagari, as
characterised in Sec.~\ref{sec:failure}.

\subsection{Quality controls}
\label{sec:controls}

Each control below was added in response to a defect observed in generated
output.

\begin{itemize}\itemsep2pt
\item \textbf{Distractor isolation.} Incorrect options are drawn from a
different chapter of the same volume, so a distractor cannot accidentally be
a true statement about the asked topic.
\item \textbf{Answer-leakage guards.} An item is rejected if its answer
appears in the stem outside the option block, or in any context field
supplied alongside it. Verdict items carry no context at all, since supplying
the source passage would reveal the verdict.
\item \textbf{OCR filtering.} Sentences are rejected if they carry running
heads, spaced-out punctuation, or an anomalous proportion of tokens absent
from an English lexicon. Verses failing the Devanagari well-formedness test
motivated in Sec.~\ref{sec:failure} are excluded from cloze and
verse-comprehension items.
\item \textbf{Option-order randomisation.} Correct answers are distributed
uniformly over option positions. An earlier build placed true/false verdicts
at fixed positions, which would have allowed a model to score on position
alone across 17\% of the set.
\end{itemize}

\subsection{Resulting distribution}

The set comprises 20{,}699 fill-in-the-blank items, 18{,}262
multiple-choice, 15{,}428 short-answer and descriptive, 11{,}430
assertion--reason, 9{,}089 multiple-selection and 5{,}092 match-the-column.
Train and validation splits are made on the fact identifier rather than the
row, so no rendering of a training fact can appear in validation.
Distribution across volumes is uneven and follows the prose/verse split of
Sec.~\ref{sec:characterisation}: Bhavaprakasha Nighantu contributes 34.4\% of
items and Sushruta Samhita 3.0\%.

% =====================================================================
\section{Base Model Selection}
\label{sec:selection}

Before fine-tuning we select a base model empirically. A 100-item evaluation
subset was drawn from the validation split --- 25 items each of
assertion--reason, multiple-choice, fill-in-the-blank and match-the-column
--- balanced across all eight volumes and disjoint from the training set.
Four instruction-tuned candidates in the 2--4B class were scored under
identical prompting.

All four task types are four-option items, so chance accuracy is 0.25 and raw
accuracy alone is misleading. We additionally report \emph{format
compliance}, the proportion of items for which the model emitted a parseable
answer, because a model that cannot follow the answer format produces a score
reflecting instruction-following rather than domain knowledge.

\begin{table}[t]
\centering
\small
\setlength{\tabcolsep}{4pt}
\begin{tabular}{@{}lrrr@{}}
\toprule
\textbf{Model} & \textbf{Acc.} & \textbf{Lift} & \textbf{Compl.} \\
\midrule
Qwen2.5-3B-Instruct     & \textbf{0.510} & +0.260 & 0.99 \\
Phi-3-mini-4k-instruct  & 0.390 & +0.140 & 0.99 \\
gemma-2-2b-it           & 0.360 & +0.110 & 1.00 \\
Llama-3.2-3B-Instruct   & 0.290 & +0.040 & 0.92 \\
\bottomrule
\end{tabular}
\caption{Base model selection over 100 held-out items. Chance accuracy is
0.25; \emph{Lift} is accuracy above chance; \emph{Compl.} is format
compliance.}
\label{tab:selection}
\end{table}

Because all models answered the same items, differences are tested with a
paired exact McNemar test rather than by comparing marginal accuracies.
Qwen2.5-3B-Instruct significantly outperforms gemma-2-2b-it ($p=0.014$) and
Llama-3.2-3B-Instruct ($p=0.0003$). Its advantage over Phi-3-mini is not
significant at $n=100$ ($p=0.058$), and the remaining three models are
mutually indistinguishable. Qwen2.5-3B-Instruct was selected.

Two observations constrain what fine-tuning must achieve. First, no model
exceeds chance on assertion--reason items, and gemma-2-2b-it falls below it
at 0.16; match-the-column is at or near chance for all four. These two
reasoning-heavy formats are where the base models hold no usable prior.
Second, two models answer by position rather than content --- Llama selects
option C on 50\% of items and gemma selects D on 41\%, against a key
distributed 24/30/20/26 --- so part of their apparent score is an artefact of
key placement. Finally, 29 of the 100 items were solved by no model, and the
union of all four reaches only 0.71, indicating substantial headroom in the
benchmark itself.

% =====================================================================
\section{Model Trimming and Pruning}
\label{sec:pruning}

Deployment targets a single local accelerator, the same constraint that
governed extraction. Vocabulary trimming and structured pruning are planned
ahead of fine-tuning, with quantised export as the deployment artefact.

\gap{Trimming and pruning --- not yet performed}{%
To complete: vocabulary trimming criterion and retained vocabulary size,
justified against the Devanagari and IAST coverage required by
Sec.~\ref{sec:characterisation}; pruning method, granularity and target
sparsity; parameter count and on-disk size before and after; and the change
in Table~\ref{tab:selection} accuracy attributable to compression alone,
measured before any fine-tuning.}

% =====================================================================
\section{Fine-Tuning Protocol}
\label{sec:finetune}

Fine-tuning proceeds from the selected base model on the instruction corpus
of Sec.~\ref{sec:synthesis}, split on the fact identifier. The evaluation
subset of Sec.~\ref{sec:selection} is held out entirely.

\gap{Fine-tuning --- not yet performed}{%
To complete: adaptation method and, if low-rank, the rank, scaling factor and
target modules; learning rate, schedule, epochs, effective batch size and
sequence length; the sampling policy over the 80{,}000 items, in particular
whether the Bhavaprakasha over-representation noted in
Sec.~\ref{sec:synthesis} was corrected; hardware and wall-clock cost; and the
checkpoint selection criterion.}

% =====================================================================
\section{Benchmarking Against AyurParam}
\label{sec:ayurparam}

The reference point for this work is AyurParam~\cite{ayurparam}, evaluated on
BhashaBench-Ayur~\cite{bhashabench}. Two comparisons are planned. The first
evaluates our fine-tuned model on our own held-out benchmark alongside
AyurParam, measuring shloka-level grounding in the eight source volumes. The
second evaluates our model on the public BhashaBench-Ayur suite, measuring
whether grounding in primary sources transfers to examination questions
written independently of our corpus.

The two directions test different claims and we report both. A model tuned on
our corpus is expected to hold an advantage on the first by construction;
only the second constitutes evidence of general domain competence.

% =====================================================================
\section{Implementation and Infrastructure}
\label{sec:infra}

\subsection{Hugging Face}

The corpus is held as a private Hugging Face dataset repository, serving
three purposes. It versions the corpus, so a change to a converter produces a
new revision rather than overwriting the artefact an experiment was run
against. It provides authenticated access, appropriate while the corpus
derives from copyrighted printed editions. And it supplies base model weights
and tokenizers to the evaluation environment through the same interface.

Gated repositories required explicit handling. Several candidates in
Sec.~\ref{sec:selection} are licence-gated, and access failures surface only
when the model configuration is fetched, producing an authorisation error
rather than an informative message. Our harness performs an access preflight
over the candidate list before any weights are downloaded, and substitutes an
equivalent ungated mirror where the gate could not be cleared.

\subsection{Kaggle execution}

Model evaluation runs on Kaggle notebook instances with a single T4
accelerator, which provides free accelerated compute adequate for models in
the 2--4B class. Three constraints of that environment shaped the harness.
Sessions are wiped on termination, so results are archived per model and
downloaded before the session closes. Sessions idle out, so long batches are
queued in groups. And accelerator memory is not released between models in a
batch, so the harness explicitly frees weights between candidates; without
this the second model in a run fails to allocate even where it would fit
alone.

The harness is distributed as a single self-contained notebook with both the
scoring engine and the evaluation items embedded, so a run requires no
attached dataset and no path configuration. Two defects found during
development are worth recording, as both silently depress scores rather than
raising an error. Decoder-only models must be left-padded for batched
generation; under default right-padding, sequences shorter than the batch
maximum have their continuation generated after a run of padding tokens and
return empty. And answer extraction must accept the response shapes models
actually produce --- a letter followed by a newline and a rationale is the
most common instruct-model output and was initially scored as a non-answer.

\subsection{Reproducibility}

Corpus analysis, question synthesis and evaluation are deterministic and
seeded; re-running the synthesis pipeline reproduces the 80{,}000-item corpus
byte-for-byte. Every generated item carries the volume, sthana, chapter,
verse identifier and chunk identifier of its source.

% =====================================================================
\section{Experiments}
\label{sec:experiments}

\gap{Experiments --- reserved}{%
To be populated once Secs.~\ref{sec:pruning}--\ref{sec:ayurparam} are
executed. Should consolidate: compression cost; fine-tuning gain per question
format with paired significance tests as in Sec.~\ref{sec:selection}; the
AyurParam comparison in both directions; and an ablation isolating which of
the quality controls in Sec.~\ref{sec:controls} materially affected
downstream accuracy.}

% =====================================================================
\section{Limitations}
\label{sec:limitations}

Three limitations are inherent to the present corpus rather than to the
pipeline. First, coverage is bounded by extraction quality: the 6.1\% of
chunks yielding no instruction items are concentrated in Sharangadhara
Samhita and Ashtanga Hridaya for the reasons given in Sec.~\ref{sec:failure},
and no downstream processing recovers content the recognition pass did not
resolve.

Second, question stems are produced from templates. Answers are verbatim
source text and therefore correct, but a model fine-tuned on templated stems
may learn the template rather than the content. Paraphrasing stems while
holding answers and option sets fixed is the natural mitigation and is not
yet applied.

Third, the instruction corpus is less shloka-centric than the source material
would suggest. Only 12.3\% of items involve the Sanskrit verse; 81.2\% rest
on prose commentary alone. This follows from the token distribution of
Sec.~\ref{sec:characterisation} --- a chunk carries one verse but may carry
many commentary sentences --- and means a model trained on the full set
learns commentary comprehension more than verse comprehension. Where verse
grounding is the objective, the set must be filtered accordingly.

Finally, the selection benchmark of Sec.~\ref{sec:selection} comprises 100
items, sufficient to separate the best candidate from the weakest but not to
separate adjacent candidates, as the non-significant Qwen--Phi-3 comparison
shows.

% =====================================================================
\section{Conclusion}

\gap{Conclusion --- reserved}{%
Write last, once Sec.~\ref{sec:experiments} carries results.}

% =====================================================================
{\small
% NOTE: no \bibliographystyle here on purpose. The references below are a
% manual thebibliography, so BibTeX never runs. Declaring ieee_fullname would
% make Overleaf's latexmk look for ieee_fullname.bst, which ships only with
% the official CVPR template and is not in a stock TeX Live installation.
\begin{thebibliography}{99}\itemsep2pt

\bibitem{surya}
Datalab.
\newblock Surya: OCR, layout analysis, reading order and table recognition in
  90+ languages.
\newblock \url{https://github.com/datalab-to/surya}.

\bibitem{marker}
V.~Paruchuri and Datalab.
\newblock Marker: convert PDF to markdown and JSON with high accuracy.
\newblock \url{https://github.com/datalab-to/marker}.

\bibitem{docling}
N.~Livathinos, C.~Auer, et~al.
\newblock Docling: an efficient open-source toolkit for AI-driven document
  conversion.
\newblock {\em arXiv:2501.17887}, 2025.

\bibitem{doclingtr}
C.~Auer, et~al.
\newblock Docling technical report.
\newblock {\em arXiv:2408.09869}, 2024.

\bibitem{ppocr}
Y.~Du, C.~Li, R.~Guo, X.~Yin, W.~Liu, J.~Zhou, Y.~Bai, Z.~Yu, Y.~Yang,
  Q.~Dang, and H.~Wang.
\newblock PP-OCR: a practical ultra lightweight OCR system.
\newblock {\em arXiv:2009.09941}, 2020.

\bibitem{ppocrv2}
Y.~Du, et~al.
\newblock PP-OCRv2: bag of tricks for ultra lightweight OCR system.
\newblock {\em arXiv:2109.03144}, 2021.

\bibitem{lfm2vl}
Liquid AI.
\newblock LFM2-VL: efficient vision-language models.
\newblock \url{https://www.liquid.ai/blog/lfm2-vl-efficient-vision-language-models}.

\bibitem{devaocr}
A comprehensive survey of OCR for Devanagari script based languages.
\newblock \TBD{complete authors/venue}.

\bibitem{devavlm}
Can OCR-VLMs read Devanagari? A stress-test benchmark and post-correction
  study.
\newblock {\em arXiv:2606.29213}.

\bibitem{sanskritdl}
A systematic framework for Sanskrit character recognition using deep learning.
\newblock {\em ELCVIA}, 24(2), 2025.

\bibitem{ayurparam}
M.~Nauman, S.~Gvm, V.~Devane, S.~Pawar, V.~Thakur, K.~Pundalik, P.~Sawarkar,
  R.~Saluja, M.~Desarkar, and G.~Ramakrishnan.
\newblock AyurParam: a state-of-the-art bilingual language model for Ayurveda.
\newblock {\em arXiv:2511.02374}, 2025.

\bibitem{bhashabench}
V.~Devane, M.~Nauman, B.~Patel, et~al.
\newblock BhashaBench V1: a comprehensive benchmark for the quadrant of Indic
  domains.
\newblock {\em arXiv:2510.25409}, 2025.

\bibitem{parambench}
ParamBench: a graduate-level benchmark for evaluating LLM understanding on
  Indic subjects.
\newblock {\em arXiv:2508.16185}, 2025.

\bibitem{indicparam}
IndicParam: benchmark to evaluate LLMs on low-resource Indic languages.
\newblock {\em arXiv:2512.00333}, 2025.

\bibitem{cpt-indic}
Adapting multilingual LLMs to low-resource languages using continued
  pre-training and synthetic corpus.
\newblock {\em arXiv:2410.14815}, 2024.

\bibitem{bhaasha}
Bhaasha, Bhasa, Zaban: a survey for low-resourced languages in South Asia ---
  current stage and challenges.
\newblock {\em arXiv:2509.11570}, 2025.

\bibitem{attrground}
From real to synthetic: synthesizing millions of diversified and complicated
  user instructions with attributed grounding.
\newblock {\em arXiv:2506.03968}, 2025.

\bibitem{cotself}
CoT-Self-Instruct: building high-quality synthetic data for reasoning and
  non-reasoning tasks.
\newblock {\em arXiv:2507.23751}, 2025.

\bibitem{seedfree}
Seed-free synthetic data generation framework for instruction-tuning LLMs: a
  case study in Thai.
\newblock {\em arXiv:2411.15484}, 2024.

\end{thebibliography}
}

\end{document}
